Kernel Search (KS) remains one of the strongest matheuristics for the
large-scale Multi-Source Capacitated Facility Location Problem (MSCFLP).
This paper takes a design-analysis perspective on G\&S~2012 Kernel Search.
A close examination of the original method reveals three design choices that
can each be refined independently, yielding consistent improvements in both
solution quality and computation time.

The first concerns LP construction cost: the full LP relaxation contains all
$|\J||\I|$ shipping variables, most of which play no role in the optimal
solution.
Column generation recovers the same bound, facility partition, and
reduced-cost signal at substantially lower cost.

The second concerns arc selection: the global reduced-cost threshold $\gamma$
used to build restricted MILPs provides no size guarantee and systematically
admits fewer arcs for bucket facilities than for kernel facilities, because
bucket facilities have LP dual zero and receive no capacity-cost discount on
their arc reduced costs.

The third concerns stage transitions: when the arc set changes between stages,
the incumbent solution from the previous subproblem may no longer be feasible
in the next, because arcs it uses may be absent.
The solver must rediscover feasibility from scratch at every stage transition.

We address each with a targeted modification.
Column generation replaces full LP construction.
A client-centric flat-$k$ criterion replaces the global $\gamma$ threshold,
selecting exactly $k$ arcs per client by reduced cost and treating all
facilities uniformly.
Stage transitions are made feasibility-preserving by including every arc used
by the current incumbent when building the next edge set; we prove that this
guarantees the incumbent is feasible in the next restricted MILP by
construction, making it a valid warm start at every stage.
The resulting repaired method retains the full architecture of G\&S~2012 and
is named \emph{Funnel Kernel Search} (FKS).

The computational study isolates each modification via a dedicated comparison,
so that each improvement can be attributed to a specific design decision.
Results are reported across five large-scale benchmark families from the
Avella et al.\ test bed.
FKS finds the proven optimal solution on 109 of the 110 Test Bed~A instances
for which an optimum is known --- solutions that the state-of-the-art Benders
decomposition of Fischetti et al.\ requires up to 50{,}000 seconds per instance
to certify --- while FKS operates within a 300-second per-subproblem budget.
On Test Bed~B and Test Bed~C, FKS establishes 22 new best-known solutions,
improving on all published upper bounds it encounters.
